\chapter{Introduction}
\label{chap:intro}

\section{Motivation}
\label{sec:intro:motivation}

\TODO{Quantum networks promise secure communication and distributed quantum
computation. Noisy Bell pairs are the central bottleneck. Purification
protocols improve fidelity at the cost of consuming additional copies.}

\section{The Scheduling Gap}
\label{sec:intro:gap}

\TODO{Bencha et al.\ established the HRGS architecture and showed that
purification can be inserted at any stage (RGSS, link, end-node) with
optimistic (blind) execution to avoid double heralding round-trips.
Their numerical results use a single hand-chosen schedule and explicitly
note that schedule optimization remains open. This thesis fills that gap.}

\section{Research Objectives and Contributions}
\label{sec:intro:contributions}

\TODO{State the three to four concrete contributions:
\begin{enumerate}
    \item A formal model of the schedule space for HRGS-based repeaters.
    \item A DAG-based cost evaluator that reproduces closed-form benchmarks
          from prior work.
    \item Exact and heuristic search algorithms that find schedules
          outperforming the baseline.
    \item A regime characterisation showing when simple schedules are
          near-optimal and when they are not.
\end{enumerate}}

\section{Thesis Outline}
\label{sec:intro:outline}

\TODO{Chapter-by-chapter roadmap.}